% Part: first-order-logic
% Chapter: models-theories
% Section: expressing-relations

\documentclass[../../../include/open-logic-section]{subfiles}

\begin{document}

\olfileid{fol}{mat}{exr}

\olsection{Relaties met formules beschrijven in \article{structure}
  \printtoken{S}{structure}}

\begin{explain}
Met !!{formula}s kun je eigenschappen en relaties in
!!a{structure}~$\Struct M$ beschrijven. Daarbij gebruik je alleen de
basistekens van de taal~$\Lang L$ van~$\Struct M$. Concreet werkt dat
zo. Het !!{domain} van $\Struct M$ is een verzameling objecten. De
!!{constant}s, !!{function}s en !!{predicate}s krijgen in~$\Struct M$
een betekenis: het zijn bepaalde objecten in~$\Domain M$, functies
op~$\Domain M$ en relaties op~$\Domain M$. Stel bijvoorbeeld dat
$\Obj A^2_0$ in $\Lang L$ staat. Dan koppelt $\Struct M$ dit teken aan
de relatie~$R = \Assign{{\Obj A^2_0}}{M}$. De formule
$\Atom{\Obj A^2_0}{\Obj v_1, \Obj v_2}$ \emph{drukt precies die
relatie uit}. Dat betekent het volgende. Stel dat een bedeling~$s$
de variabele $\Obj v_1$ koppelt aan $a \in \Domain{M}$ en $\Obj v_2$
aan $b \in \Domain M$. Dan geldt
\[
Rab \text{\quad iff\quad} \Sat{M}{\Atom{\Obj A^2_0}{\Obj v_1, \Obj v_2}}[s].
\]
De bedeling is hier echt nodig. We kunnen niet gewoon zeggen ``$Rab$
dan en slechts dan als $\Sat{M}{\Atom{\Obj A^2_0}{a, b}}$''. De dingen
$a$ en $b$ zijn namelijk geen tekens van onze taal. Het zijn
!!{element}s van~$\Domain{M}$.

We hebben niet alleen atomaire !!{formula}s. Met logische connectieven
en kwantoren kunnen we ze tot ingewikkeldere !!{formula}s combineren.
Zo'n formule kan een relatie beschrijven die niet rechtstreeks
in~$\Struct M$ is ingebouwd. We willen weten hoe dat werkt en precies
welke relaties je in !!a{structure} kunt definiëren.
\end{explain}

\begin{defn}
Neem een !!{formula} $!A(\Obj v_1,\dots, \Obj v_n)$ van $\Lang L$
waarin alleen $\Obj v_1$,\dots, $\Obj v_n$ vrij voorkomen. Neem ook
!!a{structure}~$\Struct M$ voor~$\Lang L$. We zeggen dat
$!A(\Obj v_1,\dots, \Obj v_n)$ de relatie~$R \subseteq
\Domain M^n$ \emph{uitdrukt} precies dan als
\[
Ra_1\dots a_n \text{\quad iff\quad} \Sat{M}{\Atom{!A}{\Obj
    v_1,\dots, \Obj v_n}}[s]
\]
voor iedere bedeling~$s$ waarvoor $s(\Obj v_i) = a_i$ ($i = 1,
\dots, n$). De formule is dus precies waar bij de rij objecten die in
de relatie zit.
\end{defn}

\begin{ex}
Kijk naar het standaardmodel van de rekenkunde~$\Struct N$. De
!!{formula} $\Obj v_1 < \Obj v_2 \lor \eq[\Obj v_1][\Obj v_2]$
drukt de relatie~$\le$ op~$\Nat$ uit. De !!{formula}
$\eq[\Obj v_2][\Obj v_1']$ drukt de opvolgerrelatie uit. Dat is de
relatie $R \subseteq \Nat^2$ waarvoor $Rnm$ geldt als $m$ de opvolger
van~$n$ is. De formule $\eq[\Obj v_1][\Obj v_2']$ drukt juist de
voorgangerrelatie uit. Ook kun je dezelfde relatie op meer dan één
manier beschrijven. De !!{formula}s
$\lexists[\Obj v_3][(\eq/[\Obj v_3][\Obj 0] \land
  \eq[\Obj v_2][(\Obj v_1 + \Obj v_3)])]$ en $\lexists[\Obj
  v_3][\eq[(\Obj v_1 + {\Obj v_3}')][v_2]]$ drukken allebei de relatie
$\Obj <$ uit. Het predicaatsymbool~$<$ is daarom niet echt nodig in de
taal van de rekenkunde: we kunnen het met andere tekens definiëren.
\end{ex}

\begin{explain}
Dit idee werkt niet alleen in één bepaalde !!{structure}. Het is juist
algemeen nuttig wanneer we met een taal één bedoeld model of meerdere
bedoelde modellen beschrijven, dus wanneer we een theorie geven. Zo'n
theorie gebruikt vaak maar een paar !!{predicate}s als basistekens. In
het domein waarover de theorie gaat, spelen meestal nog veel andere
relaties een rol. Kunnen we die andere relaties steeds beschrijven met
de relaties achter de basis-!!{predicate}s? Dan zeggen we dat we ze in
de taal kunnen \emph{definiëren}.
\end{explain}

\begin{prob}
Zoek !!{formula}s in $\Lang L_A$ die de volgende relaties beschrijven:
\begin{enumerate}
\item $n$ ligt tussen $i$ en $j$;
\item $n$ deelt $m$ (dus $m$ is een veelvoud van $n$);
\item $n$ is een priemgetal (dus behalve $1$ en $n$ deelt geen enkel
  getal~$n$).
\end{enumerate}
\end{prob}

\begin{prob}
Stel dat de formule $!A(\Obj v_1, \Obj v_2)$ de relatie $R \subseteq
\Domain M^2$ in !!a{structure}~$\Struct M$ uitdrukt. Zoek formules voor
de volgende relaties:
\begin{enumerate}
\item de inverse $R^{-1}$ van $R$, die elk geordend paar omdraait;
\item de samenstelling $R \mid R$, waarbij je de relatie twee keer
  achter elkaar volgt;
\end{enumerate}
Kun je ook $R^+$ uitdrukken? Dat is de transitieve afsluiting van~$R$:
de kleinste transitieve relatie die de oorspronkelijke relatie bevat.
\end{prob}

\begin{prob}
Neem een taal $\Lang{L}$ met alleen het tweeplaatsige predicaatsymbool
$<$. Er zijn dus geen andere !!{constant}s, !!{function}s of
!!{predicate}s---behalve natuurlijk~$\eq$. Neem de structuur
$\Struct{N}$ waarvoor $\Domain{N} = \Nat$ en
$\Assign{<}{N} = \Setabs{\tuple{n,m}}{n < m}$. Bewijs nu het volgende:
\begin{enumerate}
\item $\{ 0 \}$ is definieerbaar in $\Struct{N}$;
\item $\{ 1 \}$ is definieerbaar in $\Struct{N}$;
\item $\{ 2 \}$ is definieerbaar in $\Struct{N}$;
\item voor iedere $n \in \Nat$ is de verzameling $\{ n \}$
  definieerbaar in $\Struct{N}$;
\item iedere eindige deelverzameling van $\Domain{N}$ is definieerbaar
  in $\Struct{N}$;
\item iedere co-eindige deelverzameling van $\Domain{N}$ is
  definieerbaar in $\Struct{N}$. Hierbij is een $X \subseteq \Nat$
  co-eindig precies als $\Nat \setminus X$ eindig is.
\end{enumerate}
\end{prob}


\end{document}
